\section{Noise}  \label{sec:noise}


\subsection{Bad Pixels}

Bad pixels arise from cosmic rays or bad behavior of any one pixel. With CCD sensors,
pixels can become 'dead' or stuck at one value. This usually results in a bright or
dark full or partial column.

With CMOS sensors, includesion store charge. This results in a partial report
of the energy seen, or a report that includes that bit saved-back but inadvertantly
added to this readout -- resulting in a high value.

There is also column variances with CMOS sensors. The GAIN keyword should
really be a vector, one value for each column or row -- depending on the
readout direction.

One test is to take a single image, and add the columns together. Then report
the statistics on that summed array:

\begingroup \fontsize{10pt}{10pt}
\selectfont
%%\begin{Verbatim} [commandchars=\\\{\}]
\begin{verbatim} 
from astopy.io import fits
filename='file.fits'
with fits.open(filename,uint=True) as f:
   d = f[0].data
   v = d.sum(axis=0)  # the first dim = column 
print(f"variance {v.var()}, std {v.std{}, mean {v.mean()}}"
nf = fits.PrimaryHDU(data,f[0].header)
nf.writeto('colstats_'+filename,output_verify='fix',overwrite=True)
print(f"#! ds9 {'colstats_'+filename}") # handy for ds9.
\end{verbatim}
\endgroup
%% \end{Verbatim}


\subsection{Zero, Dark, Flat Consistency}

Create a data cube for a collection of files. Compute the mean, median, variance and
standard deviation across the pixels. These values can be turned into a 2D image
file to 'see' the values in relation to themselves. 

Another trick is to simply use imstat to see certain values, for example: eyeball
the variance between zeros and flats.

\subsubsection{Flat Files}

The flat cube analyis

For 11 files:

    Flat8080B_12s_20260626_085321m1.fits
    Flat8080B_12s_20260626_085321m10.fits
    Flat8080B_12s_20260626_085321m11.fits
    Flat8080B_12s_20260626_085321m2.fits
    Flat8080B_12s_20260626_085321m3.fits
    Flat8080B_12s_20260626_085321m4.fits
    Flat8080B_12s_20260626_085321m5.fits
    Flat8080B_12s_20260626_085321m6.fits
    Flat8080B_12s_20260626_085321m7.fits
    Flat8080B_12s_20260626_085321m8.fits
    Flat8080B_12s_20260626_085321m9.fits

made a cube using pyraflogin3 newcube.

Using ds9, with a view of the cube '3 1 2' \index{ds9!Cube} a tall display is seen,
where left to right is from frame 1 to 11 and the top is the blue
end at 0 and the bottom is the red end at fill size of the sensor.
The height is 6280.

Use projection to move around the image to get an intuitive feel for the
variances in the data.



\begingroup \fontsize{10pt}{10pt}
\selectfont
%%\begin{Verbatim} [commandchars=\\\{\}]
\begin{verbatim} 
from astropy.io import fits

rawdata = []
with fits.open('file.fits',uint=True) as f:
   d = f[0].data
   rawdata.append(d)

cube = np.array(rawdata)
\end{verbatim}
\endgroup
%% \end{Verbatim}
